% Part: many-valued-logics
% Chapter: sequent-calculus
% Section: structural-rules

\documentclass[../../../include/open-logic-section]{subfiles}

\begin{document}

\olfileid{mvl}{seq}{str}

\olsection{संरचनात्मक नियम}

$n$-बाजूंच्या क्रमवर्ती कलनातील संरचनात्मक नियम
अभिजात प्रसंगाप्रमाणेच कार्य करतात; मात्र ते प्रत्येक जागा~$i$
साठी स्वतंत्रपणे लागू होतात.

\begin{defish}
\begin{center}
\AxiomC{$ \Gamma_1 \nSequent \dots \nSequent \phantom{!A,}\Gamma_i \nSequent \dots \nSequent \Gamma_n $}
\RightLabel{$\iR{\Weakening}{i}$}
\UnaryInfC{$\Gamma_1 \nSequent \dots \nSequent !A, \Gamma_i \nSequent \dots \nSequent \Gamma_n$}
\DisplayProof
\\[2ex]
\AxiomC{$ \Gamma_1 \nSequent \dots \nSequent !A, !A, \Gamma_i \nSequent \dots \nSequent \Gamma_n $}
\RightLabel{$\iR{\Contraction}{i}$}
\UnaryInfC{$\Gamma_1 \nSequent \dots \nSequent \phantom{!A,}!A, \Gamma_i \nSequent
\dots \nSequent \Gamma_n$}
\DisplayProof
\\[2ex]
\AxiomC{$ \Gamma_1 \nSequent \dots \nSequent \Gamma_i, !A, !B, \Gamma_i' \nSequent \dots \nSequent \Gamma_n $}
\RightLabel{$\iR{\Exchange}{i}$}
\UnaryInfC{$\Gamma_1 \nSequent \dots \nSequent \Gamma_i, !B, !A, \Gamma_i' \nSequent \dots
\nSequent \Gamma_n$}
\DisplayProof
\end{center}
\end{defish}

दुर्बलीकरण, आकुंचन आणि अदलाबदल यांच्या सलग निगमनांची मालिका
अनेकदा दुहेरी निगमन रेषांनी दर्शवली जाते.

\Cut{} नियमाची अनेक रूपे असतात: क्रमवर्तीतील भिन्न जागांच्या
प्रत्येक जोडी~$i \neq j$ साठी एक रूप:
\begin{defish}
\[
\AxiomC{$ \Gamma_1 \nSequent \dots \nSequent !A, \Gamma_i \nSequent \dots \nSequent \Gamma_n $}
\AxiomC{$ \Delta_1 \nSequent \dots \nSequent !A, \Delta_j \nSequent \dots \nSequent \Delta_n $}
\RightLabel{$\iR{\Cut}{i,j}$}
\BinaryInfC{$\Gamma_1,\Delta_1 \nSequent \dots \nSequent \Gamma_n, \Delta_n$}
\DisplayProof
\]
\end{defish}

\end{document}
